\originalsection{从立体角到局部高斯定律}\label{tb:gauss-local}
通量是有向面积积分, 并不是一根根真实管线的数目. 场线画得多密取决于绘图约定, 而$\Phi_E=\int\ee\cdot\uv n\dd A$具有确定的单位$\mathrm{N\,m^2/C}$. 穿入面为负、穿出面为正. 对点电荷, $\ee\cdot\uv n\dd A=\kc q\dd\Omega$, 其中$\dd\Omega$是从源点看面元的有向立体角. 闭面围住源点时立体角积分为$4\pi$; 不围住时, 同一射线穿入与穿出贡献成对抵消. 这解释了高斯定律为何不要求曲面为球面.

把积分定律用于一个足够小的长方体, 对边通量之差为$\partial_xE_x\dd x\dd y\dd z$等三项. 除体积后取极限, 得
\[\nabla\cdot\ee=\rho/\eps.\]
在普通静电区域又有$\nabla\times\ee=0$, 因而$\ee=-\nabla V$给出泊松方程$\nabla^2V=-\rho/\eps$. 无体电荷区为拉普拉斯方程$\nabla^2V=0$. 它们分别规定场的源和环流; 还须给边界条件才能求得唯一的实际场. 面电荷使导数有跃变, 不应在面上假装所有偏导仍为普通光滑函数.

\begin{example}
\textbf{原4.8.2: 点电荷对正方形的通量.} \label{ex:squareflux}
边长$2\ell$的正方形位于$z=\ell$, 中心在$z$轴. 点电荷$Q$在原点, 法向取$+\uv z$. 求通量, 并用积分检验.
\end{example}
\begin{solution}
这张正方形是以原点为中心、边长$2\ell$的立方体一面. 六面受旋转对称性支配, 通量各为$Q/(6\eps)$. 原题(b)立方体总通量为$Q/\eps$. 若换成面朝原点的法向, 答案变号. 直接计算时$E_z=\kc Q\ell/(x^2+y^2+\ell^2)^{3/2}$, 故
\begin{align*}
\Phi_E&=4\kc Q\int_0^\ell\int_0^\ell
\frac{\ell\dd y\dd x}{(x^2+y^2+\ell^2)^{3/2}}\\
&=4\kc Q\int_0^\ell\frac{\ell^2\dd x}{(x^2+\ell^2)\sqrt{x^2+2\ell^2}}
=4\kc Q\left[\arctan\frac{x}{\sqrt{x^2+2\ell^2}}\right]_0^\ell
=\frac Q{6\eps}.
\end{align*}
立方体方法并未声称各面上场恒定; 它只利用六个\emph{积分}相等.
\end{solution}

\begin{example}
\textbf{原4.8.1: 两张非导体电荷面.} 
无限面位于$x=0,a$, 面密度分别为$\sigma_1,\sigma_2$. 求三区域的场. 原题是等量异号面, 以本题符号取$\sigma_1=-\sigma,\sigma_2=+\sigma$, 明确该特例.
\end{example}
\begin{solution}
一张面贡献$\sigma_i\operatorname{sgn}(x-x_i)\uv x/(2\eps)$, 因而
\[E_x=\frac1{2\eps}\begin{cases}-\sigma_1-\sigma_2,&x<0,\\ \sigma_1-\sigma_2,&0<x<a,\\ \sigma_1+\sigma_2,&x>a.\end{cases}\]
原题特例给板外零、板内$E_x=-\sigma/\eps$. 相同正密度的变式使中间零场、外部场不零; 相反等量密度则中间非零、外部零. 面电荷不能取一张面本身的奇异场代入受力公式, 下一节将明确处理.
\end{solution}

\begin{example}
\textbf{原4.8.3: 引力的高斯形式.} 
半径$a$、质量$M$的均匀球壳, 求其壳内重力场, 并写出引力的闭面定律. 另讨论实心球作为变式.
\end{example}
\begin{solution}
负号表示通量流向质量, 所以$\oint\vct g\cdot\dd\vct A=-4\pi GM_{\rm enc}$. 壳内部的高斯球不包围质量且场有球对称, 所以内场为零, 外部$-GM\uv r/r^2$. 对另一个均匀实心球变式, 球半径$a$, 总质量$M$, 内部包围质量$Mr^3/a^3$, 得$\vct g=-GMr\uv r/a^3$; 外部为$-GM\uv r/r^2$. 与电场形式相似的原因是两者都遵从平方反比叠加, 并非引力可以有正负质量随意相消.
\end{solution}

\originalsection{导体空腔、曲率与电压力}\label{tb:electric-pressure}
导体静电平衡的三个结论应按因果相连: 移动电荷停止重排, 材料内部电场为零, 故相连材料等势; 表面若有切向场仍会推动电荷, 故场垂直表面. 空腔中有电荷时内壁感应总量为负的腔内总电荷, 而不是逐点均匀分布. 原例4.5中, 净电荷$Q$的导体腔内放$q$, 内壁总量$-q$、外壁$Q+q$, 这两个结果无需腔形或$q$位置的球对称.

两相距很远的导体球通过细线连接, 半径$a,b$, 忽略彼此势的影响, 等势条件$\kc q_1/a=\kc q_2/b$给$q_1/q_2=a/b$; 表面密度之比$\sigma_1/\sigma_2=b/a$. 小球附近场较强. 这支持尖端处场集中的物理图像, 但不能把$\sigma\propto1/\text{曲率半径}$冒充任意形状导体表面的精确局部定律.

取表面小片面密度$\sigma$, 平面局部模型中小片自己的场在两侧为$\pm\sigma\uv n/(2\eps)$. 自身电荷不能对自身作净力. 其余电荷产生的场应是两侧总场的平均:
\[\ee_{\rm other}=(\ee_++\ee_-)/2.\]
导体内$\ee_-=0$, 外$\ee_+=\sigma\uv n/\eps$, 所以单位面积受到向外电压力
\[p_E=\sigma^2/(2\eps)=\eps E_+^2/2.\]
若误用$\sigma E_+$, 就把自身场也算进去了, 多出一倍. 表面电荷为负时$\sigma$与外场一起反号, 压力仍向外.

原4.6附录把场想作沿线拉紧、横向挤压的介质. 可用电应力矩阵具体表达:
\[T_{ij}=\eps(E_iE_j-\tfrac12 E^2\delta_{ij}),\qquad
\vct F=\oint T\cdot\uv n\dd A.\]
其中$\delta_{ij}$在$i=j$时为1, 否则为0; 对重复$j$从$x,y,z$三个分量求和, 矩阵散度的第$i$分量是$\sum_j\partial_jT_{ij}$.
取$\ee$沿$z$, 有$T_{zz}=+u_E$, $T_{xx}=T_{yy}=-u_E$, 分别是沿场张力和横向压力. 在静电光滑区, 用$\partial_jE_i=\partial_iE_j$展开$\partial_jT_{ij}$, 含梯度的两项抵消, 剩$\eps E_i\partial_jE_j=\rho E_i$. 积分便等于内部电荷受力. 这是教材动画中场形变与力的定量连接, 没有把场线当成物质绳索. 恒定场使有电荷粒子加速, 时间变场对固定粒子的力则随时间变; 两带电摆球在重力、电力和绳张力共同作用下达到平衡.

\originalsection{球内电势、场能与挖去电荷的叠加}\label{tb:charged-sphere}
\begin{example}
\textbf{原4.6及4.8.4: 薄壳与实心球的势.} 
薄球壳或均匀实心球半径$a$, 总电荷$Q$. 取无穷远为零势, 求内外电势.
\end{example}
\begin{solution}
两者外部都是$V=\kc Q/r$. 壳内零场故$V=\kc Q/a$. 实心球内部$E_r=\kc Qr/a^3$, 从边界积分:
\[V(r)=V(a)+\int_r^a\frac{\kc Qr'}{a^3}\dd r'
=\frac{\kc Q}{2a}\left(3-\frac{r^2}{a^2}\right),\quad r\leq a.\]
势及一阶径向导数在实心球界面连续; 二阶导数会因体密度跳变而改变. 薄壳的势连续, 场却在壳上跳变. 高斯定律定的是场, 积分常数由参考点及接续确定.
\end{solution}

\begin{example}
\textbf{原4.10.1、4.10.5的经典延伸: 球形空腔与自能.} 
均匀绝缘球体密度$\rho$, 挖去完全位于其内的球形空腔, 腔心相对大球心位移$\vct a$. 求腔内场. 另求未挖空半径$R$实心球的装配自能.
\end{example}
\begin{solution}
把“挖去”表示为叠加密度$-\rho$的小球. 腔内点相对两球心分别为$\rr,\rr-\vct a$, 故
\[\ee_{\rm cavity}=\frac\rho{3\eps}\rr-\frac\rho{3\eps}(\rr-\vct a)=\frac\rho{3\eps}\vct a.\]
场均匀且与腔半径无关, 但要求腔没有穿出大球; 这不是导体空腔屏蔽.

装配均匀球时, 半径$r$以内已有$q(r)=Q(r/R)^3$, 下一层$\dd q=3Qr^2\dd r/R^3$. 每层的外部电势在$r$为$\kc q(r)/r$, 因而
\[U=\int_0^R\frac{\kc q(r)}r\dd q=\frac{3\kc Q^2}{R^6}\int_0^Rr^4\dd r=\frac{3\kc Q^2}{5R}.\]
也可对球内、球外$\eps E^2/2$分开积分, 分别给$\kc Q^2/(10R)$和$\kc Q^2/(2R)$. 只积分球外场会漏掉内部能量.
\end{solution}

若密度只依赖$r$, 通用路线是$Q(r)=4\pi\int_0^r\rho(s)s^2\dd s$, 再用$E_r=Q(r)/(4\pi\eps r^2)$, 最后积分势. 例如$\rho=\rho_0r/R$时内场$E_r=\rho_0r^2/(4\eps R)$; 非均匀密度不能直接使用均匀球的$r^3$比例.

P--N结的入门静电模型是两片相邻无限带电耗尽层. 设$-a<x<0$密度$-\rho_-$, $0<x<b$密度$+\rho_+$, 且$\rho_-a=\rho_+b$, 外部场取零. 从$\partial_xE_x=\rho/\eps$得左层$E_x=-\rho_-(x+a)/\eps$, 右层$E_x=\rho_+(x-b)/\eps$, 在$x=0$接续. 电势用负积分求出, 总势升为$(\rho_-a^2+\rho_+b^2)/(2\eps)$. 该模型只展示空间电荷与内建场, 没有据此推出真实半导体的载流子统计或结电压数值.

\begin{center}
\begin{tikzpicture}
\begin{axis}[width=.8\linewidth,height=4.5cm,xlabel={$r/a$},ylabel={$Va/(\kc Q)$},xmin=0,xmax=3,ymin=0,ymax=1.65,samples=101]
\addplot[thick,structurecolor,domain=0:1] {(3-x^2)/2};\addplot[thick,structurecolor,domain=1:3] {1/x};
\addplot[dashed,domain=0:1] {1};\addplot[dashed,domain=1:3] {1/x};
\end{axis}
\end{tikzpicture}
\end{center}
原4.8.4的势图采用实线为均匀实心球、虚线为薄壳; 两者外部相同, 内部一个二次下降、一个恒定. 实心球中心有限且斜率零, 壳面场跃变但势连续.

分段电势还可反过来识别电荷. 球对称时
\[E_r=-\frac{\dd V}{\dd r},\qquad \rho(r)=-\eps\frac1{r^2}\frac{\dd}{\dd r}\left(r^2\frac{\dd V}{\dd r}\right).\]
若$r<R$有$V=A-Br^2$, $r>R$有$V=C/r$, 连续要求$A-BR^2=C/R$. 内部$E_r=2Br,\rho=6\eps B$, 外部$E_r=C/r^2,\rho=0$, 界面面电荷$\sigma=\eps(C/R^2-2BR)$. 若$C=2BR^3$则没有面电荷, 正是均匀体球; 否则须把表面层另计. $A,B,C$单位依次为V、V/m$^2$、V$\cdot$m, 并非任意同单位常数. 这建立原3.10.9、4.10.6逆问题的路线: 先分区导数、再用法向跃变, 不能只看外部$1/r$便把全部电荷都猜到中心.
